\section{Моноид}
	\tikzsetfigurename{LinearAlgebra_Monoid_}

\begin{definition}[Моноид]
    \emph{Моноидом} называется полугруппа с нейтральным элементом.
    Если операция является коммутативной, то можно говорить о \emph{коммутативном моноиде}.
\end{definition}

\begin{example}
    Приведём примеры моноидов, построенных на основе полугрупп из предыдущего раздела.
    \begin{itemize}
        \item Неотрицательные целые числа (или же натуральные числа с нулём) с операцией сложения являются моноидом.
              Нейтральный элемент~--- $0$.
        \item Целые числа, дополненные значением $-\infty$ (\enquote{минус бесконечность}) с операцией взятия наибольшего из двух ($\max$) являются моноидом.
              Нейтральный элемент~--- $-\infty$.
        \item Множество всех строк конечной длины с пустой строкой (строка длины 0) над фиксированным алфавитом $\Sigma$ и операцией конкатенации является моноидом.
              Нейтральный элемент~--- пустая строка.
        \item Квадратные неотрицательные матрицы%
              \sidenote{Неотрицательной называется матрица, все элементы которой не меньше нуля.} фиксированного размера с операцией умножения задают моноид.
              Нейтральный элемент~--- единичная матрица.
    \end{itemize}
\end{example}

